\section{Conclusion}
\label{sec:conclusion}

% TODO: Write 1-2 paragraphs summarizing findings.

We analyzed five real-world networks across three datasets using a consistent
pipeline of degree statistics, betweenness centrality, eigenvector centrality,
and Louvain community detection.
A recurring finding is the \textit{hub--bridge duality}: the nodes with the
highest degree (or eigenvector centrality) are not the same nodes that
maximize information flow across the network.
This has practical implications — removing a high-degree node may not
disrupt communication as much as targeting a high-betweenness bridge.

In the Higgs Twitter data, node \texttt{88} dominates in-degree and
eigenvector centrality across all three interaction types, confirming its
role as the primary information source during the Higgs boson event.
In the Enron network, the community structure reflects the underlying
organizational hierarchy, with executive inboxes acting as dense hubs
within their respective communities.
The Facebook ego-network yields 16 well-separated communities consistent
with distinct social circles.

% TODO: Add any quantitative take-aways specific to your paper's thesis.

\section{Future Work}
\label{sec:future}

Several extensions of this work are natural:

\textbf{Temporal analysis.}
The Higgs dataset includes activity timestamps in
\texttt{higgs-activity\_time.txt}.
A natural follow-up is to study how centrality rankings evolve over time and
whether early adopters of the Higgs boson story maintain their influence.

\textbf{Multiplex centrality.}
Because mention, reply, and retweet networks share the same node set, one
could compute \textit{multiplex centrality} measures that aggregate
information across layers, identifying nodes that are simultaneously
important across all interaction types.

\textbf{Robustness and targeted removal.}
Given the hub--bridge distinction, a comparative simulation of random vs.\
targeted node removal (targeting hubs vs.\ bridges) would quantify the
resilience of each network.

\textbf{Feature-based community interpretation.}
The Facebook dataset includes rich node feature vectors (\texttt{node\_to\_features.txt}).
Linking detected communities to demographic or interest features could
validate the semantic meaning of the Louvain partitions.
